Đối với bài toán phân lớp nhị phân, mô hình xây dựng một hàm tuyến tính:
\begin{equation}
  a(\boldsymbol{x}) = \boldsymbol{w}^{\intercal} \boldsymbol{x} + w_0,
\end{equation}

sau đó ánh xạ giá trị này thông qua hàm sigmoid: \begin{equation}
  \sigma(a) = \frac{1}{1 + e^{-a}}.
\end{equation}

Khi đó đầu ra của mô hình được viết là: \begin{equation}
  y = \sigma(\boldsymbol{w}^{\intercal} \boldsymbol{x} + w_0).
\end{equation}

Đối với bài toán phân lớp đa lớp, mô hình sử dụng một vector điểm số (logits)
cho từng lớp: \begin{equation}
  a_k(\boldsymbol{x}) = \boldsymbol{w}_k^{\intercal} \boldsymbol{x} + w_{k0},
  \quad k = 1, \dots, K.
\end{equation}

Các giá trị này được chuẩn hóa thông qua hàm softmax: \begin{equation}
  y_k = \frac{\exp(a_k)}{\sum_{j=1}^{K} \exp(a_j)}.
\end{equation}

Hàm mất mát được sử dụng trong huấn luyện là hàm cross-entropy:
\begin{equation}
  E(\boldsymbol{W}) = - \sum_{n=1}^{N} \sum_{k=1}^{K}t_{nk} \log y_{nk},
\end{equation} trong đó $t_{nk}$ là nhãn one-hot và $y_{nk}$ là đầu ra của mô
hình.

Gradient của hàm mất mát theo tham số $\boldsymbol{w}_k$ được tính như sau:
\begin{equation}
  \nabla_{\boldsymbol{w}_k} E =
  \sum_{n=1}^{N} (y_{nk} - t_{nk}) \boldsymbol{x}_n,
\end{equation} trong đó $y_{nk}$ là đầu ra tại lớp $k$ của mẫu thứ $n$.


Mô hình Logistic Regression có thể được xem như một mô hình tuyến tính kết hợp
với một hàm kích hoạt phi tuyến, cho phép xây dựng ranh giới quyết định phức
tạp hơn trong không gian đặc trưng thông qua tối ưu hóa hàm mất mát lồi.
